\section{Uso de ferramentas de inteligência artificial}
\label{sec:ia}

Usamos o Claude (Anthropic), pelo Claude Code no VS Code, como apoio
durante todo o projeto:

\begin{itemize}[nosep]
  \item \textbf{Código:} escrever a primeira versão dos programas a
  partir da especificação que o grupo preparou (pipeline, constantes
  comentadas, ferramenta de gabarito e script de avaliação).
  \item \textbf{Testes e diagnóstico:} gerar vídeos sintéticos para
  testar o código e analisar os nossos vídeos e gravações da webcam.
  Exemplos: identificar a exposição automática como causa do fundo que
  sumia e a borda larga dos gabaritos.
  \item \textbf{Correções:} propor e testar a contagem regressiva, as
  versões da compensação de exposição, a trava da webcam e os ajustes
  de borda.
  \item \textbf{Relatório:} organizar as anotações de cada teste, gerar
  as figuras a partir dos nossos resultados e redigir um rascunho deste
  texto, que o grupo revisou.
\end{itemize}

As gravações, os testes com a webcam, a pintura dos gabaritos, a escolha
dos casos difíceis e as decisões sobre o que manter foram feitos pelo
grupo, que revisou o código e o texto e é responsável por explicá-los.
